\Needspace{6\baselineskip}
\begin{center}\hypertarget{native-250-article178}{}\pdfbookmark[2]{178.}{article-178}178.\end{center}
Die soeben gefundene Funktion kann die Wahrscheinlichkeiten der Fehler gewiss nicht mit aller Strenge ausdrücken. Da die möglichen Fehler stets auf bestimmte Grenzen beschränkt sind, müsste die Wahrscheinlichkeit größerer Fehler stets $=0$ werden, während unsere Formel immer einen endlichen Wert ergibt. Dieser Mangel, der jeder analytischen Funktion ihrer Natur nach anhaften muss, ist für alle praktischen Anwendungen ohne Belang. Der Wert unserer Funktion nimmt so rasch ab, sobald $h\Delta$ einen beträchtlichen Wert erreicht hat, dass man ihn ohne Bedenken als gleichwertig mit $0$ ansehen kann. Außerdem lässt die Natur der Sache es niemals zu, die Grenzen der Fehler selbst mit absoluter Strenge anzugeben.

\SourceLeaf{251}Die Konstante $h$ lässt sich ferner als Maß für die Genauigkeit der Beobachtungen ansehen. Die Wahrscheinlichkeit eines Fehlers $\Delta$ werde in einer Beobachtungsreihe durch $\frac{h}{\BareRoot{\pi}}e^{-hh\Delta\Delta}$ ausgedrückt, in einer anderen Reihe mehr oder weniger genauer Beobachtungen dagegen durch $\frac{h'}{\BareRoot{\pi}}e^{-h'h'\Delta\Delta}$. Die Erwartung, dass der Fehler bei einer Beobachtung der ersten Reihe zwischen den Grenzen $-\delta$ und $+\delta$ liegt, wird durch das Integral $\int\frac{h}{\BareRoot{\pi}}e^{-hh\Delta\Delta}\,\mathrm{d}\Delta$ von $\Delta=-\delta$ bis $\Delta=+\delta$ ausgedrückt. Ebenso wird die Erwartung, dass der Fehler bei einer Beobachtung der zweiten Reihe innerhalb der Grenzen $-\delta'$ und $+\delta'$ bleibt, durch das Integral $\int\frac{h'}{\BareRoot{\pi}}e^{-h'h'\Delta\Delta}\,\mathrm{d}\Delta$ von $\Delta=-\delta'$ bis $\Delta=+\delta'$ ausgedrückt. Beide Integrale sind offenbar gleich, wenn $h\delta=h'\delta'$ gilt. Ist also zum Beispiel $h'=2h$, so kann in der ersten Reihe ein doppelt so großer Fehler ebenso leicht auftreten wie ein halb so großer in der zweiten. Im gewöhnlichen Sprachgebrauch schreibt man den Beobachtungen der zweiten Reihe dann die doppelte Genauigkeit zu.
